\documentclass[11pt,a4paper]{article}
\usepackage[utf8]{inputenc}
\usepackage[T1]{fontenc}
\usepackage[spanish,es-nodecimaldot]{babel}
\usepackage{lmodern}
\usepackage{amsmath,amssymb}
\usepackage{booktabs}
\usepackage{geometry}
\usepackage{xcolor}
\usepackage{hyperref}
\usepackage{microtype}
\geometry{margin=2.4cm}
\hypersetup{colorlinks=true,linkcolor=blue!55!black,urlcolor=blue!55!black}
\newcommand{\vect}[1]{\boldsymbol{#1}}
\newcommand{\atan}{\operatorname{atan2}}

\title{Desarrollo matemático de la cinemática y la marcha cartesiana\\
Nova Spot Micro}
\author{Proyecto de tesis --- robot cuadrúpedo Nova Spot Micro}
\date{13 de agosto de 2026}

\begin{document}
\maketitle
\begin{abstract}
Este documento registra la formulación matemática implementada para convertir
posiciones cartesianas de los pies en referencias articulares del modelo
Nova Spot Micro y generar una marcha de gateo estable en Gazebo. Se presentan
las convenciones, la cinemática directa, la solución analítica de cinemática
inversa, la trayectoria por fases, las restricciones y la validación realizada.
Las dimensiones dinámicas y articulares continúan siendo provisionales hasta
medir el robot físico.
\end{abstract}
\tableofcontents

\section{Objetivo y alcance}
El controlador inicial utilizaba listas fijas de ángulos. Aunque una versión
conservadora mantenía el equilibrio, el robot casi no avanzaba. La modificación
consistió en definir el movimiento en el espacio cartesiano de cada pie y
calcular en cada muestra los tres ángulos de su pata mediante cinemática
inversa. La política implementada sigue siendo control convencional nominal;
no corresponde todavía a la capa de aprendizaje por refuerzo de la tesis.

La formulación toma como referencia conceptual el proyecto MIT
\href{https://github.com/mike4192/spot_micro_kinematics_python}{spot\_micro\_kinematics\_python}
de Mike4192, pero fue rederivada para las longitudes, ejes ROS y signos del URDF
propio. No se trasladaron directamente sus dimensiones ni su numeración.

\section{Convenciones y parámetros geométricos}
El marco local de una cadera coincide, antes de las rotaciones, con el marco del
cuerpo: $x$ apunta hacia delante, $y$ hacia la izquierda y $z$ hacia arriba.
Cada pata tiene tres variables:
\[
  \vect q_i = [q_{c,i},\ q_{f,i},\ q_{t,i}]^T,
\]
donde los subíndices indican coxa, fémur y tibia. Se usa
$s=+1$ para patas izquierdas y $s=-1$ para patas derechas. El eje de la coxa
en el URDF es $s\vect e_x$; fémur y tibia rotan alrededor del eje $y$ local.

\begin{table}[h]
\centering
\caption{Parámetros geométricos provisionales empleados.}
\begin{tabular}{lll}
\toprule
Símbolo & Valor & Significado\\
\midrule
$l_c$ & $0.090$ m & longitud de coxa\\
$l_f$ & $0.105$ m & longitud de fémur\\
$l_t$ & $0.132$ m & longitud de tibia\\
$(q_c,q_f,q_t)_0$ & $(0.10,0.42,-0.84)$ rad & postura nominal\\
\bottomrule
\end{tabular}
\end{table}

Los límites provisionales usados como supervisor cinemático son
\[
-0.60\le q_c\le0.60,\qquad
-1.20\le q_f\le1.20,\qquad
-2.20\le q_t\le0.10.
\]

\section{Cinemática directa de una pata}
Primero se resuelve el plano fémur--tibia. La posición horizontal y vertical
antes de aplicar la rotación de coxa es
\begin{align}
x_p &= -l_f\sin q_f-l_t\sin(q_f+q_t),\\
z_p &= -l_f\cos q_f-l_t\cos(q_f+q_t).
\end{align}
El desplazamiento lateral sin rotación es $y_p=s l_c$. La coxa gira el vector
$(y_p,z_p)$ un ángulo $a=sq_c$ alrededor de $x$. Por tanto,
\begin{align}
x &= x_p,\\
y &= \cos(a)\,s l_c-\sin(a)\,z_p,\\
z &= \sin(a)\,s l_c+\cos(a)\,z_p.
\end{align}
La función completa queda
\[
\boxed{
\vect p_F(q_c,q_f,q_t,s)=
\begin{bmatrix}
-l_f\sin q_f-l_t\sin(q_f+q_t)\\
\cos(sq_c)s l_c-\sin(sq_c)z_p\\
\sin(sq_c)s l_c+\cos(sq_c)z_p
\end{bmatrix}}
\]
y se implementa como \texttt{forward\_leg}.

\section{Cinemática inversa}
Se desea encontrar $\vect q$ a partir de un objetivo
$\vect p_F=[x,y,z]^T$.

\subsection{Separación de la coxa}
La rotación de coxa conserva la norma en el plano $yz$. Por geometría,
\[
r_{yz}^2=y^2+z^2=l_c^2+z_p^2.
\]
Se elige la rama de la pata que apunta hacia abajo:
\[
\boxed{z_p=-\sqrt{y^2+z^2-l_c^2}}.
\]
Una condición necesaria de alcanzabilidad es
$y^2+z^2-l_c^2\ge0$. El ángulo aplicado alrededor de $x$ se obtiene como
la diferencia entre el argumento del vector final y el inicial:
\begin{align}
a &= \atan(z,y)-\atan(z_p,s l_c),\\
q_c &= \boxed{\frac{1}{s}\atan(\sin a,\cos a)}.
\end{align}
La segunda expresión normaliza $a$ al intervalo principal y respeta el signo
del eje de las patas derecha e izquierda.

\subsection{Solución del plano fémur--tibia}
Se definen coordenadas positivas convenientes
\[
d=-z_p,\qquad b=-x.
\]
Por ley de cosenos,
\[
D=\frac{d^2+b^2-l_f^2-l_t^2}{2l_fl_t}.
\]
Debe cumplirse $-1\le D\le1$. Se selecciona la configuración de rodilla
flexionada compatible con el URDF:
\[
\boxed{q_t=-\arccos D}.
\]
Finalmente, la descomposición geométrica del manipulador plano da
\[
\boxed{
q_f=\atan(b,d)-
\atan\!\left(l_t\sin q_t,\ l_f+l_t\cos q_t\right)}.
\]
Antes de aceptar la solución se comprueba que los tres valores sean finitos y
se encuentren en los límites articulares. Si falla alguna condición, la marcha
se rechaza y el controlador conserva la postura.

\section{Comprobación FK--IK}
La prueba automatizada toma la postura nominal para ambos valores de $s$:
\[
\vect q_0=(0.10,0.42,-0.84)\ \text{rad}.
\]
Calcula $\vect p_0=FK(\vect q_0,s)$ y luego
$\widehat{\vect q}_0=IK(\vect p_0,s)$. La condición verificada es
\[
|q_{0,j}-\widehat q_{0,j}|<10^{-9}\quad\forall j.
\]
También se genera un ciclo completo y se comprueba que el mayor salto entre
muestras consecutivas sea menor que $0.20$ rad. Las dos pruebas aprobaron.

\section{Trayectoria cartesiana de gateo}
Cada pata parte de su posición neutral $\vect p_{0,i}$ obtenida por FK. La fase
global discretizada es
\[
\phi_k=\frac{k}{N},\qquad k=0,\ldots,N-1.
\]
Para cada pata se aplica un desfase $\delta_i$ y un módulo unitario:
\[
\phi_i=(\phi_k+\delta_i)\bmod 1.
\]
Los desfases implementados producen el orden de oscilación
delantera izquierda, trasera derecha, delantera derecha y trasera izquierda.
El factor de apoyo es $\beta=0.75$, por lo que solamente una pata entra en
oscilación a la vez.

\subsection{Fase de apoyo}
Para $0\le\phi_i<\beta$ se define
\[
u=\frac{\phi_i}{\beta},\qquad
x_i=x_{0,i}+L\left(\frac12-u\right),\qquad z_i=z_{0,i}.
\]
El pie se desplaza hacia atrás respecto del cuerpo, representando el avance del
cuerpo sobre un apoyo aproximadamente fijo.

\subsection{Fase de oscilación}
Para $\beta\le\phi_i<1$ se usa
\[
v=\frac{\phi_i-\beta}{1-\beta},
\]
\begin{align}
x_i &= x_{0,i}+L\left(-\frac12+v\right),\\
z_i &= z_{0,i}+4H v(1-v).
\end{align}
La parábola vertical vale cero en los extremos y alcanza exactamente $H$ en
$v=1/2$. La coordenada lateral se mantiene como $y_i=y_{0,i}$. Cada punto
$(x_i,y_i,z_i)$ se convierte en ángulos mediante IK.

\section{Discretización y temporización}
La línea base validada utiliza
\[
N=24,\quad L=0.018\ \text{m},\quad H=0.014\ \text{m},
\quad \Delta t=0.18\ \text{s}.
\]
Por ello, el periodo nominal es
\[
\boxed{T=N\Delta t=24(0.18)=4.32\ \text{s}}.
\]
La trayectoria articular usa un tiempo de transición del $90\%$ de cada fase:
\[
t_{\mathrm{tr}}=0.90\Delta t=0.162\ \text{s}.
\]
El resto de la fase ofrece un margen pequeño antes de enviar el punto siguiente.

\section{Restricciones paramétricas}
El controlador valida los parámetros antes de iniciar el gateo:
\begin{align*}
8&\le N\le80, & N&\equiv0\pmod4,\\
0.002&\le L\le0.040\ \text{m},\\
0.004&\le H\le0.030\ \text{m},\\
0.08&\le\Delta t\le1.0\ \text{s}.
\end{align*}
Estas cotas son de simulación y no autorizan su transferencia automática a los
MG996R. La calibración física debe establecer límites definitivos.

\section{Validación en Gazebo}
Se tomaron doce muestras separadas aproximadamente cuatro segundos, cubriendo
más de diez ciclos. Los resultados fueron:
\begin{table}[h]
\centering
\begin{tabular}{ll}
\toprule
Métrica & Resultado\\
\midrule
Posición longitudinal inicial & $-0.00212$ m\\
Posición longitudinal final & $0.21676$ m\\
Desplazamiento & $0.21887$ m\\
Altura mínima--máxima & $0.22287$--$0.22428$ m\\
Desviación lateral absoluta máxima & $0.00986$ m\\
Componente $|w|$ mínima del cuaternión & $0.99906$\\
Caídas & 0\\
\bottomrule
\end{tabular}
\end{table}

Tomando aproximadamente $44$ s entre la primera y la última muestra, la
velocidad media estimada es
\[
\bar v_x\approx\frac{0.21887}{44}=4.97\times10^{-3}\ \text{m/s}.
\]
La variación total observada de altura fue
\[
\Delta z=0.22428-0.22287=1.41\times10^{-3}\ \text{m}.
\]
La cota angular global asociada únicamente a $w$ puede expresarse como
$\theta=2\arccos|w|$; con $|w|=0.99906$ corresponde aproximadamente a
$0.087$ rad, es decir, $5.0^{\circ}$. Esta magnitud no sustituye el cálculo
separado de roll, pitch y yaw. El nodo de métricas implementado posteriormente
ya realiza esa conversión y la publica junto con los estados articulares.

\section{Limitaciones y siguiente desarrollo}
El resultado demuestra avance estable en el modelo actual, no valida todavía
un gemelo digital ni el robot físico. Permanecen provisionales masas, inercias,
fricción, esfuerzos, ganancias y geometría. La marcha no desplaza explícitamente
el cuerpo para maximizar margen estático; su estabilidad se observó en la
configuración actual de Gazebo.

Los contactos de los cuatro pies y el supervisor provisional de altura e
inclinación ya están integrados en Gazebo. Un ensayo de 76 ciclos mostró solo
32.550\% de coincidencia con el plan ideal y ausencia de despegue trasero. Los
siguientes pasos son corregir esa trayectoria, repetir la medición, calcular el
polígono de soporte con contactos válidos, proyectar el centro de masa y definir
el margen de estabilidad. Después se podrá comparar sistemáticamente la
marcha nominal y la misma marcha con correcciones acotadas de aprendizaje por
refuerzo.

\end{document}
